\input{../tex-inputs/latex-korrekturansicht-vorspann}

\section[Ida Fulda u. a. an Arthur Schnitzler, 8. 12. 1901]{L04347 Ida Fulda u. a. an Arthur Schnitzler, 8. 12. 1901}
\nopagebreak\mylabel{L04347v}
\rehead{ }\normalsize\beginnumbering\briefempfaengerindex{Schnitzler, Arthur@\textsc{Schnitzler, Arthur}!zzz, Christine@\emph{von Christine }!1901-12-081@{{[}8. 12. 1901{]}}|(be}\briefempfaengerindex{Schnitzler, Arthur@\textsc{Schnitzler, Arthur}!zzzFriedmann, Emma@\emph{von Emma Friedmann}!1901-12-081@{{[}8. 12. 1901{]}}|(be}\briefempfaengerindex{Schnitzler, Arthur@\textsc{Schnitzler, Arthur}!zzzFriedmann, Alfred@\emph{von Alfred Friedmann}!1901-12-081@{{[}8. 12. 1901{]}}|(be}\briefempfaengerindex{Schnitzler, Arthur@\textsc{Schnitzler, Arthur}!zzzGoldmann, Paul@\emph{von Paul Goldmann}!1901-12-081@{{[}8. 12. 1901{]}}|(be}\briefempfaengerindex{Schnitzler, Arthur@\textsc{Schnitzler, Arthur}!zzzOlden, Hans@\emph{von Hans Olden}!1901-12-081@{{[}8. 12. 1901{]}}|(be}\briefempfaengerindex{Schnitzler, Arthur@\textsc{Schnitzler, Arthur}!zzzd’Albert, Ida@\emph{von Ida d’Albert}!1901-12-081@{{[}8. 12. 1901{]}}|(be}
\toendnotes[C]{\smallbreak\pagebreak[2]}
\correspDesc{Versand  durch Ida d’Albert, Hans Olden, Paul Goldmann, Alfred Friedmann, Emma Friedmann, Christine am [8. 12. 1901] in Berlin
\newline{}Übermittlung  am 8. 12. 1901 in Berlin
\newline{}Erhalt  durch Arthur Schnitzler am 10. 12. 1901 in Wien}\toendnotes[C]{\smallbreak}
\Standort{CUL, Schnitzler, B 331.}
\physDesc{Postkarte, 385 Zeichen
\newline{}Handschrift Ida d’Albert: blaue Tinte, lateinische Kurrent
\newline{}Handschrift Hans Olden: blaue Tinte, deutsche Kurrent
\newline{}Handschrift Paul Goldmann: blaue Tinte
\newline{}Handschrift Alfred Friedmann: blaue Tinte, deutsche Kurrent
\newline{}Handschrift Emma Friedmann: blaue Tinte, lateinische Kurrent
\newline{}Handschrift Christine : blaue Tinte, deutsche Kurrent
\newline{}Versand: 1) Stempel: »\nobreak{}\oindex{Berlin@\textbf{Berlin}, \emph{Hauptstadt}|pwk}Berlin W, 8. 12. 1901, 6–7N\nobreak{}«.   2) Stempel: »\nobreak{}\oindex{IX., Alsergrund@\textbf{IX., Alsergrund}, \emph{Verwaltungsgebiet}|pwk}Wien \textcolor{gray}{9/3} 72, 10. 12. 1901, 8.V, Bestellt\nobreak{}«. 
\newline{}Schnitzler: mit Bleistift datiert: »8. 12. 90\textcolor{gray}{1}« }\toendnotes[C]{\smallbreak}\pstart{}{\pb}Herrn\pend{}\pstart{}Dr. Arthur Schnitzler\pend{}\pstart{}\textcolor{pink}{Wien}\oindex{Wien@\textbf{Wien}, \emph{Verwaltungsgebiet}|pw}{}\ledrightnote{\textcolor{pink}{Wien}}\pend{}\pstart{}\textcolor{pink}{IX. Frangaſse 1}\oindex{Wien@\textbf{Wien}!IX., Alsergrund@\textbf{IX., Alsergrund}!Frankgasse 1@\textbf{Frankgasse 1}, \emph{Wohngebäude}|pw}{}\ledrightnote{\textcolor{pink}{Frankgasse 1}}.\pend{}{\bigskip}\vspace{1em}
\pstart
           \noindent{}{\pb}Wir haben soeben mit Begeisterung Ihr
            Wohl und das der »Literatenseelen« ausgebracht, was wir Ihnen nicht vorenthalten wollen.\pend
           
\pstart
           Herzliche Grüsse{\\[\baselineskip]}\spacefill\mbox{Idus Fulda}\pend
           \leftskip=0em{}\selectlanguage{ngerman}\vspace{1em}\pstart {[}hs. Christine:{]} eine \spacefill\mbox{Christine}\pend{}
\pstart
           \noindent{}{[}hs. d’Albert:{]} und was für Eine!\pend
           \selectlanguage{ngerman}\vspace{1em}\pstart \spacefill\mbox{{[}hs. Goldmann:{]} Paul Goldmann.}\pend{}\selectlanguage{ngerman}\vspace{1em}\pstart {[}hs. Friedmann:{]} Die Sie sehr verehrende \label{K_L04347-1v}\edtext{Wirtin}{\lemma{\textnormal{\emph{Wirtin}}}\Cendnote{\textnormal{\textcolor{blue}{Emma Friedmann}\pwindex{Friedmann, Emma 1861 – 24.\,8.\,1922 Frankfurt am Main@\textsc{Friedmann, Emma} (1861 – 24.\,8.\,1922 Frankfurt am Main)|pwk}, Paul Goldmann an Arthur Schnitzler, 13. 12. [1901].}}}\label{K_L04347-1}.\pend{}\selectlanguage{ngerman}\vspace{1em}\pstart {[}hs. Olden:{]} Herzlichen Gruß von Ihrem \spacefill\mbox{Hans Olden}\pend{}\selectlanguage{ngerman}\vspace{1em}\pstart {[}hs. Friedmann:{]} Schließlich auch von dem Ihnen unbekannten Wirth \label{K_L04347-2v}\edtext{von’s Janze}{\lemma{\textnormal{\emph{von’s Janze}}}\Cendnote{\textnormal{dialektal: vom Ganzen}}}\label{K_L04347-2}: \spacefill\mbox{D\textsuperscript{r} AlfredFriedmann}\pend{}\selectlanguage{ngerman}\endnumbering\briefempfaengerindex{Schnitzler, Arthur@\textsc{Schnitzler, Arthur}!zzz, Christine@\emph{von Christine }!1901-12-081@{{[}8. 12. 1901{]}}|)be}\briefempfaengerindex{Schnitzler, Arthur@\textsc{Schnitzler, Arthur}!zzzFriedmann, Emma@\emph{von Emma Friedmann}!1901-12-081@{{[}8. 12. 1901{]}}|)be}\briefempfaengerindex{Schnitzler, Arthur@\textsc{Schnitzler, Arthur}!zzzFriedmann, Alfred@\emph{von Alfred Friedmann}!1901-12-081@{{[}8. 12. 1901{]}}|)be}\briefempfaengerindex{Schnitzler, Arthur@\textsc{Schnitzler, Arthur}!zzzGoldmann, Paul@\emph{von Paul Goldmann}!1901-12-081@{{[}8. 12. 1901{]}}|)be}\briefempfaengerindex{Schnitzler, Arthur@\textsc{Schnitzler, Arthur}!zzzOlden, Hans@\emph{von Hans Olden}!1901-12-081@{{[}8. 12. 1901{]}}|)be}\briefempfaengerindex{Schnitzler, Arthur@\textsc{Schnitzler, Arthur}!zzzd’Albert, Ida@\emph{von Ida d’Albert}!1901-12-081@{{[}8. 12. 1901{]}}|)be}\mylabel{L04347h}
\begin{anhang}
\end{anhang}\newcommand{\dateiname}{L04347}\newcommand{\titel}{Ida Fulda u. a. an Arthur Schnitzler, 8. 12. 1901}\newcommand{\editorInnen}{Herausgegeben von Jahnke, SelmaMüller, Martin Anton}\input{../tex-inputs/latex-korrekturansicht-abspann}
